%=============================================================================
% CHERN--SIMONS EIGENVALUE SPECTRUM ON CP^2 x S^3
% Complete specification of the weights w_k and eigenvalues lambda_k
% that produce beta_{U(1)} = 3.7969
% Extracted from: Unified Theory v108 (Appendix K), Ab Initio paper (v12)
% Generated 2026-03-22
%=============================================================================
\documentclass[12pt,a4paper]{article}
\usepackage{amsmath,amssymb}
\usepackage{booktabs}
\usepackage{geometry}
\geometry{margin=2.5cm}
\usepackage{hyperref}
\usepackage{longtable}

\title{The Chern--Simons Eigenvalue Spectrum on $\mathbb{C}P^2 \times S^3$\\[6pt]
\large Explicit Weights, Eigenvalues, and the Sum Producing
$\beta_{U(1)} = 3.7969$}
\author{Extracted from G.\ Alcock, \textit{Density Field Dynamics}}
\date{March 2026}

\begin{document}
\maketitle

%=============================================================================
\section{Origin and Definition}
\label{sec:origin}
%=============================================================================

The DFD internal geometry is $X = \mathbb{C}P^2 \times S^3$.  The $S^3$
factor carries an $SU(2)_k$ Chern--Simons theory at each level $k$.
The partition function of $SU(2)_k$ Chern--Simons theory on $S^3$ is the
exact result (Witten, 1989):
\begin{equation}
\label{eq:ZCS}
Z_{SU(2)_k}(S^3) = \sqrt{\frac{2}{k+2}}\;\sin\!\left(\frac{\pi}{k+2}\right).
\end{equation}

The shift $k \to k_{\mathrm{eff}} \equiv k + 2$ arises from the dual Coxeter
number $h^\vee = 2$ for $SU(2)$, required for modular invariance, quantum
consistency, and proper WZW central charge $c = 3k/(k+2)$.

%=============================================================================
\section{The Weight Function}
\label{sec:weights}
%=============================================================================

The Euclidean microsector weight at CS level $k$ is defined as the
squared modulus of the partition function:
\begin{equation}
\label{eq:weight}
\boxed{
w(k) \;=\; \bigl|Z_{SU(2)_k}(S^3)\bigr|^2
     \;=\; \frac{2}{k+2}\;\sin^2\!\left(\frac{\pi}{k+2}\right).
}
\end{equation}

This is Eq.~(5) of the Ab Initio paper and appears in Appendix~K of the
Unified Theory.

\paragraph{Physical interpretation.}
\begin{itemize}
\item The factor $2/(k+2)$ reflects the $1/k$ scaling of the effective
  coupling $g^2 \sim 1/k$: low-$k$ sectors are strongly quantum (``loud''),
  high-$k$ sectors are weakly coupled (``quiet'').
\item The $\sin^2(\pi/(k+2))$ factor encodes the SU(2) representation
  theory structure of the CS path integral on $S^3$.
\item Together, $w(k)$ peaks sharply at $k = 0$ and falls as
  $\sim 2\pi^2/(k+2)^3$ for large $k$.
\end{itemize}

%=============================================================================
\section{The Eigenvalues (Shifted Levels)}
\label{sec:eigenvalues}
%=============================================================================

The ``CS eigenvalue'' at level $k$ is the shifted level:
\begin{equation}
\label{eq:eigenvalue}
\lambda_k \;\equiv\; k_{\mathrm{eff}} \;=\; k + 2.
\end{equation}

This is the natural quantum variable in $SU(2)$ Chern--Simons theory.
The shift by $h^\vee = 2$ means the physical spectrum runs from
$\lambda_0 = 2$ (at $k = 0$) to $\lambda_{59} = 61$ (at $k = 59$).

\paragraph{Connection to $S^3$ heat kernel.}
The heat kernel on $S^3$ has eigenvalues $\lambda_n^{(\mathrm{Lapl})}
= n(n+2)/R^2$ with degeneracy $(n+1)^2$ (Appendix~K).
The shifted CS level $k+2$ matches the natural spectral parameter of the
$S^3$ Laplacian, connecting the CS microsector to the spectral geometry
of the internal space.

%=============================================================================
\section{The UV Cutoff: $k_{\max} = 60$}
\label{sec:cutoff}
%=============================================================================

The sum is truncated at $k_{\max} = 60$, derived from the Bridge Lemma
(a closed Spin$^c$ index on $\mathbb{C}P^2$):
\begin{equation}
\label{eq:kmax}
k_{\max} = \chi\!\left(\mathbb{C}P^2,\;\mathcal{O}(9) \oplus
\mathcal{O}^{\oplus 5}\right)
= \binom{11}{2} + 5 = 55 + 5 = 60.
\end{equation}

The twist bundle $E = \mathcal{O}(9) \oplus \mathcal{O}^{\oplus 5}$ is
fixed by two independent requirements:
\begin{itemize}
\item $\mathcal{O}(9) = \mathcal{O}(3)^{\otimes 3}$: minimal globally
  well-defined hypercharge twist, with $q_1 = 3$ from anomaly cancellation.
\item $\mathcal{O}^{\oplus 5}$: one factor per chiral multiplet type per SM
  generation $\{Q_L, u_R, d_R, L_L, e_R\}$.
\end{itemize}

The sum therefore runs over $k = 0, 1, 2, \ldots, 59$ (60 levels total,
standard SU(2) WZW/CS convention).

%=============================================================================
\section{The CS Weighted Average: $\beta_{U(1)} = 3.7969$}
\label{sec:beta}
%=============================================================================

The $U(1)$ lattice coupling is identified with the vacuum expectation value
of the shifted CS level:
\begin{equation}
\label{eq:beta-def}
\boxed{
\beta_{U(1)} \;=\; \langle k_{\mathrm{eff}} \rangle
\;=\; \frac{\displaystyle\sum_{k=0}^{k_{\max}-1}
  \lambda_k \cdot w(k)}
  {\displaystyle\sum_{k=0}^{k_{\max}-1} w(k)}
\;=\; \frac{\displaystyle\sum_{k=0}^{59}
  (k+2)\;\frac{2}{k+2}\sin^2\!\frac{\pi}{k+2}}
  {\displaystyle\sum_{k=0}^{59}
  \frac{2}{k+2}\sin^2\!\frac{\pi}{k+2}}
}
\end{equation}

Simplifying the numerator: $(k+2) \cdot w(k) = 2\sin^2(\pi/(k+2))$, so:
\begin{equation}
\beta_{U(1)} = \frac{\displaystyle\sum_{k=0}^{59}
  2\sin^2\!\frac{\pi}{k+2}}
  {\displaystyle\sum_{k=0}^{59}
  \frac{2}{k+2}\sin^2\!\frac{\pi}{k+2}}
= \frac{\displaystyle\sum_{k=0}^{59}
  \sin^2\!\frac{\pi}{k+2}}
  {\displaystyle\sum_{k=0}^{59}
  \frac{1}{k+2}\sin^2\!\frac{\pi}{k+2}}.
\end{equation}

%=============================================================================
\section{Explicit Numerical Evaluation}
\label{sec:numerics}
%=============================================================================

\subsection{Partial Sums}

The numerator and denominator evaluate to:
\begin{align}
\sum_{k=0}^{59} (k+2)\,w(k) &= \sum_{k=0}^{59} 2\sin^2\!\frac{\pi}{k+2}
= 8.32449\ldots \label{eq:num}\\[4pt]
\sum_{k=0}^{59} w(k) &= \sum_{k=0}^{59} \frac{2}{k+2}\sin^2\!\frac{\pi}{k+2}
= 2.19242\ldots \label{eq:den}
\end{align}

Therefore:
\begin{equation}
\boxed{
\beta_{U(1)} = \frac{8.32449}{2.19242} = 3.796945\ldots \approx 3.80.
}
\end{equation}

\subsection{Python Verification}

\begin{verbatim}
import math
def w(k): return (2.0/(k+2)) * math.sin(math.pi/(k+2))**2
Z = sum(w(k) for k in range(60))          # = 2.19242...
N = sum((k+2)*w(k) for k in range(60))    # = 8.32449...
beta = N / Z                               # = 3.796945...
\end{verbatim}

%=============================================================================
\section{Complete Eigenvalue Table}
\label{sec:table}
%=============================================================================

Table~\ref{tab:spectrum} lists all 60 CS levels with their weights,
eigenvalues, and contributions to the weighted average.  The sum is
overwhelmingly dominated by the first $\sim 10$ levels; the remaining
50 levels contribute less than 5\% of the total weight.

\begin{longtable}{rrrr}
\caption{Complete Chern--Simons eigenvalue spectrum on
$\mathbb{C}P^2 \times S^3$ at $k_{\max} = 60$.
Columns: CS level $k$; shifted level (eigenvalue)
$\lambda_k = k+2$; weight $w(k) = \frac{2}{k+2}\sin^2\frac{\pi}{k+2}$;
weighted eigenvalue $\lambda_k \cdot w(k) = 2\sin^2\frac{\pi}{k+2}$.}
\label{tab:spectrum}\\
\toprule
$k$ & $\lambda_k = k{+}2$ & $w(k)$ & $\lambda_k \cdot w(k)$ \\
\midrule
\endfirsthead
\multicolumn{4}{c}{\textit{(continued from previous page)}}\\
\toprule
$k$ & $\lambda_k = k{+}2$ & $w(k)$ & $\lambda_k \cdot w(k)$ \\
\midrule
\endhead
\midrule
\multicolumn{4}{r}{\textit{(continued on next page)}}\\
\endfoot
\bottomrule
\endlastfoot
 0 &  2 & 1.000000 & 2.000000 \\
 1 &  3 & 0.500000 & 1.500000 \\
 2 &  4 & 0.250000 & 1.000000 \\
 3 &  5 & 0.138197 & 0.690983 \\
 4 &  6 & 0.083333 & 0.500000 \\
 5 &  7 & 0.053787 & 0.376510 \\
 6 &  8 & 0.036612 & 0.292893 \\
 7 &  9 & 0.025995 & 0.233956 \\
 8 & 10 & 0.019098 & 0.190983 \\
 9 & 11 & 0.014431 & 0.158746 \\
10 & 12 & 0.011165 & 0.133975 \\
11 & 13 & 0.008811 & 0.114544 \\
12 & 14 & 0.007074 & 0.099031 \\
13 & 15 & 0.005764 & 0.086455 \\
14 & 16 & 0.004758 & 0.076120 \\
15 & 17 & 0.003972 & 0.067528 \\
16 & 18 & 0.003350 & 0.060307 \\
17 & 19 & 0.002852 & 0.054183 \\
18 & 20 & 0.002447 & 0.048943 \\
19 & 21 & 0.002116 & 0.044427 \\
20 & 22 & 0.001841 & 0.040507 \\
21 & 23 & 0.001612 & 0.037083 \\
22 & 24 & 0.001420 & 0.034074 \\
23 & 25 & 0.001257 & 0.031417 \\
24 & 26 & 0.001118 & 0.029058 \\
25 & 27 & 0.000998 & 0.026955 \\
26 & 28 & 0.000895 & 0.025072 \\
27 & 29 & 0.000806 & 0.023379 \\
28 & 30 & 0.000728 & 0.021852 \\
29 & 31 & 0.000660 & 0.020470 \\
30 & 32 & 0.000600 & 0.019215 \\
31 & 33 & 0.000548 & 0.018071 \\
32 & 34 & 0.000501 & 0.017027 \\
33 & 35 & 0.000459 & 0.016070 \\
34 & 36 & 0.000422 & 0.015192 \\
35 & 37 & 0.000389 & 0.014384 \\
36 & 38 & 0.000359 & 0.013639 \\
37 & 39 & 0.000332 & 0.012950 \\
38 & 40 & 0.000308 & 0.012312 \\
39 & 41 & 0.000286 & 0.011720 \\
40 & 42 & 0.000266 & 0.011169 \\
41 & 43 & 0.000248 & 0.010657 \\
42 & 44 & 0.000231 & 0.010179 \\
43 & 45 & 0.000216 & 0.009732 \\
44 & 46 & 0.000202 & 0.009314 \\
45 & 47 & 0.000190 & 0.008923 \\
46 & 48 & 0.000178 & 0.008555 \\
47 & 49 & 0.000168 & 0.008210 \\
48 & 50 & 0.000158 & 0.007885 \\
49 & 51 & 0.000149 & 0.007579 \\
50 & 52 & 0.000140 & 0.007291 \\
51 & 53 & 0.000132 & 0.007019 \\
52 & 54 & 0.000125 & 0.006762 \\
53 & 55 & 0.000119 & 0.006518 \\
54 & 56 & 0.000112 & 0.006288 \\
55 & 57 & 0.000106 & 0.006069 \\
56 & 58 & 0.000101 & 0.005862 \\
57 & 59 & 0.000096 & 0.005665 \\
58 & 60 & 0.000091 & 0.005478 \\
59 & 61 & 0.000087 & 0.005300 \\
\midrule
\multicolumn{2}{r}{\textbf{Totals:}} & $\mathbf{2.19242}$ & $\mathbf{8.32449}$ \\
\end{longtable}

%=============================================================================
\section{Cumulative Weight Analysis}
\label{sec:cumulative}
%=============================================================================

To show how the sum converges:

\begin{table}[h]
\centering
\caption{Cumulative fraction of total weight as a function of
included levels.}
\label{tab:cumulative}
\begin{tabular}{rcc}
\toprule
Levels included & Cumulative $\sum w(k)$ & Fraction of total \\
\midrule
$k = 0$ only     & 1.0000 & 45.6\% \\
$k = 0$--$1$     & 1.5000 & 68.4\% \\
$k = 0$--$2$     & 1.7500 & 79.8\% \\
$k = 0$--$4$     & 1.9714 & 89.9\% \\
$k = 0$--$9$     & 2.1036 & 95.9\% \\
$k = 0$--$19$    & 2.1653 & 98.8\% \\
$k = 0$--$59$    & 2.1924 & 100\% \\
\bottomrule
\end{tabular}
\end{table}

The $k = 0$ level alone contributes 45.6\% of the total weight.
The first 10 levels ($k = 0$--$9$) account for 95.9\%.  This
justifies the physical interpretation: the vacuum stiffness that
sets $\alpha$ is dominated by the strongly quantum low-$k$ CS
modes.

%=============================================================================
\section{Sensitivity to UV Cutoff}
\label{sec:sensitivity}
%=============================================================================

The dependence of $\beta_{U(1)}$ on the truncation point $k_{\max}$:

\begin{table}[h]
\centering
\caption{$\beta_{U(1)}$ vs.\ UV cutoff $k_{\max}$.  Only
$k_{\max} = 60$ (from the Bridge Lemma) yields $\alpha = 1/137$.}
\label{tab:kmax-scan}
\begin{tabular}{rcc}
\toprule
$k_{\max}$ & $\beta_{U(1)} = \langle k{+}2 \rangle$ & $\alpha$ result \\
\midrule
50 & 3.77 & $1/137$ ($+1.3\%$) \\
\textbf{60} & \textbf{3.80} & $\mathbf{1/137}$ ($\mathbf{+0.5\%}$) \\
$\infty$ & 3.94 & $1/303$ ($-55\%$, \textbf{ruled out}) \\
\bottomrule
\end{tabular}
\end{table}

%=============================================================================
\section{How $\beta_{U(1)}$ Enters the $\alpha$ Derivation}
\label{sec:chain}
%=============================================================================

The value $\beta_{U(1)} = 3.7969$ is one of eight Table~XXVIII components
in the derivation of $\alpha^{-1} = 137.036$.  It enters through two
independent routes:

\paragraph{Route A (Lattice, $\sim 1\%$).}
$\beta_{U(1)} = 3.80$ is the bare $U(1)$ lattice coupling.  Combined with
$\beta_{SU(2)} = 6 \times 3.80 = 22.80$ (Wilson ratio from stiffness ratio
$\times$ generation count), lattice Monte Carlo measures the renormalized
stiffness $\kappa_{U(1)} \approx 7.27$, giving:
\begin{equation}
\alpha^{-1} = 6\pi\,\kappa_{U(1)} \approx 137.0 \pm 1.8.
\end{equation}

\paragraph{Route B (Spectral action, sub-ppm).}
The spectral action on Toeplitz-truncated $\mathbb{C}P^2 \times S^3$ at
$d = k_{\max} + 4 = 64$ computes $\kappa_{U(1)}$ analytically from the
$a_4$ Seeley--DeWitt coefficient, incorporating all Table~XXVIII components
including the traceless projection $(d{-}1)/d = 63/64$ and the BOOST factor
$d^2/(d^2{-}1) = 4096/4095$, yielding:
\begin{equation}
\alpha^{-1} = 6\pi \times 7.26999\ldots = 137.03599985
\quad (-0.006\text{ ppm from experiment}).
\end{equation}

%=============================================================================
\section{Summary}
\label{sec:summary}
%=============================================================================

\begin{enumerate}
\item The CS eigenvalue spectrum on $\mathbb{C}P^2 \times S^3$ consists
  of 60 levels ($k = 0, \ldots, 59$) with shifted eigenvalues
  $\lambda_k = k + 2$ running from 2 to 61.

\item The weight at each level is $w(k) = \frac{2}{k+2}\sin^2\frac{\pi}{k+2}$,
  derived from the squared $SU(2)_k$ CS partition function on $S^3$.

\item The CS weighted average is:
\[
\beta_{U(1)} = \frac{\sum_{k=0}^{59}(k+2)\,w(k)}{\sum_{k=0}^{59}w(k)}
= \frac{8.32449}{2.19242} = 3.796945\ldots
\]

\item The UV cutoff $k_{\max} = 60$ is derived from the Bridge Lemma:
  $k_{\max} = \chi(\mathbb{C}P^2, \mathcal{O}(9) \oplus \mathcal{O}^{\oplus 5}) = 60$.

\item No adjustable parameters enter this computation.  Every ingredient
  ($w(k)$, $\lambda_k$, $k_{\max}$) is fixed by geometry, topology, or
  SM content.

\item The sum is dominated by low-$k$ modes: the first 10 levels carry
  95.9\% of the total weight.
\end{enumerate}

%=============================================================================
\section*{Source References}
%=============================================================================

\begin{itemize}
\item \textbf{Unified Theory v108}, Appendix~K: Microsector physics,
  Bridge Lemma, weight function definition, $k_{\max}$ discovery scan.

\item \textbf{Ab Initio paper} (v12, December 2025): Eqs.~(3)--(5),
  (14)--(16); Tables~I, IV, V.
  \url{https://doi.org/10.5281/zenodo.19173548}

\item \textbf{Alpha Inverse Formula}: Complete derivation chain
  (Table~XXVIII components), Routes A and B.
\end{itemize}

\end{document}
